% !PRINTED 16
\weberpage{16}
%<<CONT
number $\alpha$ can be produced by very different numerical sequences.
These numerical sequences, however, are all to be regarded as equal to
one another. Among other things, one may take all numbers of $S$ to be
rational fractions. Conversely, it is also easily shown that for every
given number $\alpha$ one can always give numerical sequences $S$ by
which $\alpha$ is produced, so that the totality of the numerical
sequences $S$ likewise forms a continuous set.

\begin{center}
\rule{0.18\linewidth}{0.4pt}
\end{center}

In the explanation of the fundamental operations an inconvenient
restriction has appeared in subtraction, from which we can free
ourselves by introducing \gesperrt{zero and the negative numbers}.

Let $x$ denote any element of the number system defined so far, which
we shall now call the system of \gesperrt{positive numbers}. We take
this number system a second time and, to distinguish it, denote in this
second system, which is to be called the system of \gesperrt{negative
numbers}, every element by $-x$. This second system we now order in
exactly the reverse way from the first, so that wherever in the system
$x$ ``greater'' stands, in the system $-x$ ``smaller'' is put, and
conversely. Addition and subtraction in $-x$ are explained just as in
$x$, so that $(-x)+(-y)=-(x+y)$, $(-x)-(-y)=-(x-y)$ shall hold.

But we wish to order these two number systems together in such a way
that every $-x$ is to be smaller than every $x$. We thus obtain an
ordered set in which there is no greatest and no least element. This
system is also in general continuous; only the single cut $(-x,x)$ is
produced by no element, and here, where the two systems meet, a
violation of continuity is still present. To restore continuity we
must therefore add, corresponding to the cut $(-x,x)$, a number
\gesperrt{zero, or 0}, which is precisely defined by this cut. Then
we have an ordered continuous set unbounded on both sides, the complete
series of \gesperrt{real numbers}.

In this enlarged number domain we now explain addition generally by
putting, by definition:
%CONT>>
